problem:  
Let \( (M^m,g) \) be a closed Riemannian manifold with bounded geometry and \( (N^n,h) \) a compact target with nonpositive curvature.  
For a sequence of \( \varepsilon \)-approximately harmonic maps \(u_\varepsilon : M \to N\) satisfying
\[
\|\tau(u_\varepsilon)\|_{L^2(M)} \le \varepsilon, \qquad E(u_\varepsilon) \le E_0,
\]
denote by \(\mu_\varepsilon = |\nabla u_\varepsilon|^2\,d\mathrm{vol}_g\) the energy measures, and suppose \( \mu_\varepsilon \rightharpoonup \mu = |\nabla u_*|^2\,d\mathrm{vol}_g + \nu \) for some weak limit \(u_*:M\to N\).  

Classically, the *defect measure* \( \nu \) captures the lack of strong \(W^{1,2}\)-convergence due to bubble formation. For genuine harmonic maps between manifolds, the bubbling phenomenon is well understood: each point of concentration arises via energy quantization in terms of entire harmonic spheres, and \(\nu\) corresponds to the sum of their energies.  

**Problem:**  
Assume \( (N,h) \) has nonpositive curvature and \( (M,g) \) has bounded geometry. Does there exist, for each dimension \(m\), a *universal quantitative rigidity principle* of the following type:

> There exists a dimension-dependent modulus \( \Psi_m:(0,\infty)\to(0,\infty)\), with \(\Psi_m(t)\to 0\) as \(t\to 0\), such that whenever  
> \[
> \|\tau(u_\varepsilon)\|_{L^2(M)} \le t,
> \]
> one has the **quantitative defect bound**
> \[
> \nu(M) \le \Psi_m(t)\,E_0.
> \]
In other words, can the total bubbling (defect measure) for approximate harmonic maps be *quantitatively controlled* purely in terms of the smallness of the tension field and the total energy, with constants depending only on the dimension—independent of the target geometry within the class of bounded-geometry nonpositively curved manifolds?  

Equivalently, does vanishing of the \(L^2\)-norm of the tension enforce *dimensionally uniform partial rigidity* that quantitatively precludes large-scale energy concentration?  

Why is it a "good" problem:  
This question reformulates the heart of the compactness theory for approximate harmonic maps in explicitly quantitative form, aiming at a *universal scaling law* linking analytic smallness of the tension to geometric suppression of bubbling.  
It would unify the qualitative defect–measure compactness theory with quantitative stability analysis, providing a “soft” metric for how far an approximate harmonic map can be from genuine harmonicity before energy concentration becomes inevitable.  
Unlike known almost-monotonicity or local ε-regularity results (which are local and depend on curvature constants), this problem asks for a *dimension-only* global quantitative rigidity mechanism across the entire bounded-geometry class, connecting fine-scale PDE analysis to geometric measure theory of concentration sets.  
A positive answer would introduce a new invariant \(\Psi_m\) reflecting dimensionally intrinsic stability of the harmonic map equation and might yield analogous quantitative rigidity theories for other geometric variational problems (minimal immersions, Yang–Mills fields, etc.).  
Its statement is simple, yet solving it would require deep new connections between nonlinear PDE regularity, geometric stability, and concentration–compactness analysis—fully matching the ideal of a "narrow entrance, profound depth" problem.